\documentclass[11pt,a4paper]{article}

\usepackage{fontspec}
\usepackage[english,polish]{babel}
\usepackage[a4paper,margin=25mm,headheight=22pt]{geometry}
\usepackage{microtype}
\usepackage{xcolor}
\usepackage{booktabs}
\usepackage{longtable}
\usepackage{array}
\usepackage{enumitem}
\usepackage{fancyhdr}
\usepackage{hyperref}
\usepackage{fancyvrb}

\setmainfont{TeX Gyre Pagella}
\setsansfont{TeX Gyre Heros}
\setmonofont{Latin Modern Mono}[Scale=0.88]
\definecolor{Ink}{HTML}{211E1A}
\definecolor{Accent}{HTML}{765F49}
\definecolor{Paper}{HTML}{FCFAF6}
\definecolor{Rule}{HTML}{D8CFC4}
\pagecolor{Paper}
\color{Ink}
\hypersetup{colorlinks=true,linkcolor=Accent,urlcolor=Accent,pdftitle={PoemTeX User Guide}}
\setlength{\parindent}{0pt}
\setlength{\parskip}{0.55em}
\setlist{itemsep=0.2em,topsep=0.35em}
\pagestyle{fancy}
\fancyhf{}
\fancyhead[L]{\sffamily\small PoemTeX}
\fancyhead[R]{\sffamily\small User guide}
\fancyfoot[C]{\sffamily\small\thepage}
\renewcommand{\headrulewidth}{0.3pt}
\renewcommand{\headrule}{\hbox to\headwidth{\color{Rule}\leaders\hrule height \headrulewidth\hfill}}
\newcommand{\code}[1]{\texttt{#1}}
\newcommand{\directive}[1]{\texttt{@#1}}
\newcommand{\yes}{\textcolor{Accent}{\textbullet}\ }
\newenvironment{poemcode}{\VerbatimEnvironment\begin{Verbatim}[fontsize=\small,frame=single,rulecolor=\color{Rule},xleftmargin=1em,xrightmargin=1em]}{\end{Verbatim}}
\newcolumntype{L}[1]{>{\raggedright\arraybackslash}p{#1}}

\begin{document}
\selectlanguage{english}
\hypersetup{pageanchor=false}

\begin{titlepage}
  \vspace*{0.19\textheight}
  {\sffamily\footnotesize\color{Accent}\MakeUppercase{Write poems, not TeX}\par}
  \vspace{1.2em}
  {\fontsize{34}{39}\selectfont PoemTeX\par}
  \vspace{0.25em}
  {\Large User guide\par}
  \vspace{2.5em}
  {\large A small language for writing poems and building beautiful PDFs\par}
  \vfill
  {\small Version 0.2.0 \quad Linux \quad \texttt{.poem} format 0.1\par}
\end{titlepage}
\hypersetup{pageanchor=true}

\tableofcontents
\clearpage

\section{The essentials: what you need to know}

PoemTeX lets you write plain \code{.poem} files. You do not write or edit TeX:
the \code{poem} command validates your text, generates hidden LaTeX, and builds
the PDF. A collection has one main file, normally \code{book.poem}, and any
number of included poem fragments. Only the collection is a finished book;
individual fragments can still be previewed.

\subsection{Start a collection}

From this repository, use \code{./poem}. If you installed the package, use
\code{poem} instead. First check the machine, create a project, and enter it:

\begin{poemcode}
./poem doctor
./poem init my-poems
cd my-poems
\end{poemcode}

The new project has this simple shape:

\begin{poemcode}
my-poems/
  book.poem              main collection file
  poems/pierwszy.poem    one poem fragment
  theme/theme.toml       optional visual customization
\end{poemcode}

\subsection{Write and include poems}

The main \code{book.poem} contains book information and includes, not verse:

\begin{poemcode}
@collection
@title "Światło i cisza"
@author "Anna Kowalska"
@language pl
@theme literary-classic
@paper a5

@section "Światło"
@include poems/poranek.poem
@include poems/okno.poem
\end{poemcode}

Each included file contains one poem:

\begin{poemcode}
@title "Poranek"

Pierwsza linia,
    linia z wcięciem.

Druga strofa ma *cichy akcent*
i **mocniejsze słowo**.
\end{poemcode}

Create and include a new fragment automatically with:

\begin{poemcode}
./poem new "Nocne światło"
\end{poemcode}

\subsection{Check, preview, and build}

Run these commands inside the project directory:

\begin{longtable}{@{}L{0.29\textwidth}L{0.65\textwidth}@{}}
\toprule
\textbf{Command} & \textbf{Purpose} \\
\midrule
\endhead
\code{./poem check} & Validate every source without running TeX. \\
\code{./poem preview} & Build, open the PDF, and rebuild whenever a source changes. Press Ctrl+C to stop. \\
\code{./poem build} & Produce the finished PDF in \code{build/}. \\
\code{./poem gallery} & Build comparable previews of all eight themes. \\
\bottomrule
\end{longtable}

To focus on one poem while retaining its collection's language, paper, author,
and custom theme, run \code{./poem preview poems/poranek.poem}. A fragment cannot
be built as a finished collection by itself.

\subsection{The five writing rules worth remembering}

\begin{enumerate}
  \item One nonblank source line is one poetic line. Blank lines separate stanzas.
  \item Leading and repeated spaces are preserved; trailing spaces are ignored.
  \item Use \code{*italic*}, \code{**bold**}, and \code{\{smallcaps: Text\}}.
  \item Use \code{@align left}, \code{@align center}, or \code{@align right}, then
        close that block with \code{@end}. Use \code{@space 1} through \code{@space 5}
        for a deliberate larger gap.
  \item Never paste raw TeX into a poem. TeX characters are treated as safe text,
        and files generated under \code{build/} should not be edited.
\end{enumerate}

That is enough for everyday use. The remainder is a complete reference.

\clearpage
\section{Complete reference}

\subsection{Requirements and installation}

PoemTeX supports Linux and Python 3.11 or newer. PDF output requires an ordinary
TeX Live installation with LuaLaTeX and the common packages \code{babel},
\code{fontspec}, \code{geometry}, \code{microtype}, \code{needspace},
\code{fancyhdr}, \code{hyperref}, and \code{xcolor}. \code{latexmk} is
recommended; without it PoemTeX runs LuaLaTeX twice. \code{xdg-open} is optional
and opens previews. Run \code{./poem doctor} for an exact local report.

The repository launcher needs no Python installation:

\begin{poemcode}
./poem --help
\end{poemcode}

For a system-independent development installation, create a virtual environment:

\begin{poemcode}
python3 -m venv .venv
.venv/bin/pip install -e .
.venv/bin/poem --help
\end{poemcode}

\subsection{File model, encoding, and safety}

All sources use UTF-8 and the \code{.poem} extension. A byte-order mark is
accepted. The first meaningful line of a collection must be \directive{collection};
a file without it is a fragment. A collection may contain directives only.
Verse belongs in fragments and enters the book through explicit \directive{include}
directives, in written order.

Include paths are relative to the directory containing the main collection.
They must end in \code{.poem}, may not escape the project directory, may not point
to another collection, and may occur only once. Theme files have the same
project-boundary protection. No language construct can inject raw TeX; special
characters such as \code{\%}, \code{\#}, \code{\&}, braces, dollar signs, and
backslashes are escaped by the compiler.

\subsection{Collection directives}

Quoted and unquoted values are accepted. Quoting is useful when a value contains
spaces. Directive names are case-insensitive, though lowercase is conventional.

\begin{longtable}{@{}L{0.29\textwidth}L{0.47\textwidth}L{0.16\textwidth}@{}}
\toprule
\textbf{Directive} & \textbf{Meaning and accepted values} & \textbf{Default} \\
\midrule
\endhead
\directive{collection} & Marks the main file; first meaningful line, no value. & required \\
\directive{title TEXT} & Collection title. & required \\
\directive{author TEXT} & Collection author. & required \\
\directive{subtitle TEXT} & Optional subtitle. & empty \\
\directive{language pl} & \code{pl}/\code{polish} or \code{en}/\code{english}. & \code{pl} \\
\directive{theme NAME} & Select a built-in theme. & \code{literary-classic} \\
\directive{theme-file PATH} & Load safe TOML overrides from inside the project. & none \\
\directive{paper a5} & \code{a5}, \code{a4}, or \code{letter}. & \code{a5} \\
\directive{dedication TEXT} & Front-matter dedication. & empty \\
\directive{epigraph TEXT} & Front-matter epigraph. & empty \\
\directive{epigraph-author TEXT} & Attribution; normally used with \directive{epigraph}. & empty \\
\directive{contents yes} & \code{yes/true/on} or \code{no/false/off}. & \code{yes} \\
\directive{page-numbers yes} & Enable or disable printed page numbers. & \code{yes} \\
\directive{smart-quotes yes} & Convert balanced straight double quotes typographically. & \code{yes} \\
\directive{poem-pages new} & \code{new} starts each poem on a page; \code{flow} does not. & \code{new} \\
\directive{section TEXT} & Insert a named collection division at this position. & optional \\
\directive{include PATH} & Insert one poem fragment at this position. & repeatable \\
\bottomrule
\end{longtable}

A collection-level author is inherited for a focused fragment preview. A poem
may provide its own author when needed.

\subsection{Poem headers, body directives, and comments}

Before its first body line, a fragment accepts \directive{title} (required),
\directive{subtitle}, \directive{author}, \directive{dedication},
\directive{epigraph}, and \directive{epigraph-author}. Each header directive
may appear at most once. After the body starts, metadata directives are no
longer accepted.

Alignment blocks override the theme's default body alignment:

\begin{poemcode}
@align right
tekst po prawej
jeszcze jedna linia
@end

@space 3
kolejna część
\end{poemcode}

The value of \directive{align} is \code{left}, \code{center}, or \code{right}.
Blocks cannot nest and every block must end with a value-less \directive{end}.
\directive{space N} accepts a whole number from 1 to 5 and creates a controlled
vertical break; when it follows an ordinary blank stanza break, it replaces it.
A full-line private comment begins with \code{;;} after optional indentation.

\subsection{Whitespace, wrapping, and inline text}

Whitespace is part of the public format:

\begin{itemize}
  \item Every nonblank body line creates one poetic line in order.
  \item Leading spaces and repeated spaces within a line are visually preserved.
  \item Tabs advance to four-space tab stops and trigger a warning; use spaces.
  \item Trailing whitespace has no effect. Any run of blank lines becomes one
        stanza break, so accidental extra blank lines do not change pagination.
  \item Long lines wrap at word boundaries with the theme's continuation indent.
        Automatic hyphenation is disabled. An unbroken token longer than 55
        characters triggers a possible-overflow warning.
  \item PoemTeX strongly discourages page breaks within a stanza, but an unusually
        long stanza may flow to the next page instead of overflowing.
\end{itemize}

Inline forms are intentionally few: \code{*text*} is emphasis (normally italic),
\code{**text**} is strong text (normally bold), and \code{\{smallcaps: text\}}
uses small capitals. A backslash escapes \code{*}, a backslash, or a straight
double quote. Unpaired styling markers are printed literally. With Polish smart
quotes enabled, balanced \code{"..."} pairs become Polish quotation marks;
apostrophes and unmatched quotes are not guessed or rewritten.

\subsection{Built-in themes and temporary comparison}

\begin{longtable}{@{}L{0.25\textwidth}L{0.69\textwidth}@{}}
\toprule
\textbf{Theme} & \textbf{Character} \\
\midrule
\endhead
\code{literary-classic} & Warm, restrained general-purpose literary book; the default. \\
\code{modern-minimal} & Airy contemporary sans-serif editorial design. \\
\code{quiet-manuscript} & Intimate monospaced body with restrained serif headings. \\
\code{romantic-garden} & Ivory ground, rose details, and rounded lyrical type. \\
\code{midnight-ink} & Dark navy, screen-oriented atmospheric edition. \\
\code{editorial-serif} & Confident journal typography with crisp sans-serif headings. \\
\code{spare-haiku} & Centered headings, left-set verse, and radical whitespace. \\
\code{bold-broadsheet} & Large poster-like display titles and compact serif text. \\
\bottomrule
\end{longtable}

List themes with \code{./poem themes}; add \code{--details} for measurements and
fonts. Compare without editing \code{book.poem} using
\code{./poem preview --theme modern-minimal} or build a separate result using
\code{./poem build --theme quiet-manuscript -o build/manuscript}.

\subsection{Custom themes: the metatemplate}

Themes are safe design-token files, not TeX. Print a complete starting file with:

\begin{poemcode}
./poem themes --example romantic-garden
\end{poemcode}

Save its output as, for example, \code{theme/theme.toml}, then add
\code{@theme-file theme/theme.toml} to the collection. The first setting chooses
a built-in foundation:

\begin{poemcode}
extends = "literary-classic"
body_size = 11.2
line_height = 16.0
stanza_gap_em = 1.15
accent_color = "6F5944"
background_color = "FFFCF8"
\end{poemcode}

Every supported token is listed below. Sizes and line heights are points;
margins are millimetres; gaps and continuation indent are multiples of the
current type size; colors are six hexadecimal digits without \code{\#}.

\begin{longtable}{@{}L{0.34\textwidth}L{0.60\textwidth}@{}}
\toprule
\textbf{Tokens} & \textbf{What they control} \\
\midrule
\endhead
\code{body\_font}, \code{title\_font} & Installed font-family names. \\
\code{body\_size}, \code{line\_height} & Verse size and baseline distance. \\
\code{poem\_title\_size}, \code{collection\_title\_size} & Display type sizes. \\
\code{margin\_inner\_mm}, \code{margin\_outer\_mm}, \code{margin\_top\_mm}, \code{margin\_bottom\_mm} & Page margins. \\
\code{stanza\_gap\_em}, \code{poem\_title\_gap\_em}, \code{section\_gap\_em} & Vertical rhythm. \\
\code{continuation\_indent\_em} & Hanging indent for wrapped poetic lines. \\
\code{title\_align}, \code{poem\_align}, \code{section\_align} & \code{left}, \code{center}, or \code{right}. \\
\code{page\_number\_position} & \code{outer}, \code{center}, or \code{hidden}. \\
\code{ornament} & Short plain-text divider, at most 12 characters. \\
\code{text\_color}, \code{accent\_color}, \code{background\_color} & Text, detail, and paper colors. \\
\code{title\_letter\_spacing}, \code{section\_letter\_spacing} & Tracking for display text. \\
\code{collection\_title\_case}, \code{poem\_title\_case}, \code{section\_case} & \code{as-written}, \code{uppercase}, or \code{lowercase}. \\
\code{title\_page\_top\_fraction} & Vertical title-page starting position. \\
\code{section\_title\_size} & Section heading size. \\
\bottomrule
\end{longtable}

Unknown tokens, invalid choices, unsafe font names, and out-of-range measurements
are errors. PoemTeX also rejects body leading below 110\% of the body size and
text/background contrast below 3:1. Run \code{./poem check} after each design
change. The exact numeric ranges and current defaults can be inspected with
\code{./poem themes --example THEME}.

\subsection{Command-line reference}

\begin{longtable}{@{}L{0.37\textwidth}L{0.57\textwidth}@{}}
\toprule
\textbf{Form} & \textbf{Behavior} \\
\midrule
\endhead
\code{poem init NAME} & Create a complete starter project. The target must be absent or empty. \\
\code{poem new TITLE} & Create a slug-named poem under \code{poems/} and append its include. \\
\code{poem new TITLE --filename X.poem} & Choose a simple file name explicitly. \\
\code{poem new TITLE --directory DIR} & Choose an in-project fragment directory. \\
\code{poem new TITLE --no-include} & Create the fragment without changing the collection. \\
\code{poem new TITLE --collection ROOT} & Select a non-default main file. \\
\code{poem check [SOURCE]} & Parse and validate, but do not invoke TeX. \\
\code{poem build [SOURCE] [-o DIR]} & Build a collection PDF. \\
\code{poem build --theme NAME} & Temporarily override the built-in theme. \\
\code{poem preview [SOURCE]} & Build, open, and continuously watch all dependencies. \\
\code{poem preview --no-open} & Watch without launching a viewer. \\
\code{poem preview --once} & Build a preview once and exit. \\
\code{poem gallery [SOURCE] [-o DIR]} & Build identical comparisons of every theme. \\
\code{poem gallery --open} & Also open the generated HTML index. \\
\code{poem themes [--details]} & List built-in themes and optionally their principal tokens. \\
\code{poem themes --example NAME} & Print a complete valid TOML metatemplate. \\
\code{poem doctor} & Check Python, TeX commands, packages, and theme fonts. \\
\code{poem --version} & Print the installed PoemTeX version. \\
\bottomrule
\end{longtable}

When no source is given, PoemTeX prefers \code{book.poem}; otherwise it uses the
only collection in the current directory. If several collections exist, name
one explicitly. The default output directory is \code{build/}. Generated
\code{.tex}, logs, and supporting files are implementation output: inspect them
only for diagnosis and never maintain them as source.

\subsection{Diagnostics and troubleshooting}

Errors identify the source file and line when possible and usually offer a fix.
Invalid UTF-8, unknown or duplicate directives, missing required metadata,
unsafe or repeated includes, nested/unclosed alignment blocks, and invalid
theme values stop the build. Tabs, empty poems, orphaned epigraph attributions,
very long tokens, and empty sections are warnings.

During continuous preview, an invalid edit does not destroy the last successful
PDF: PoemTeX reports the error and keeps watching. Common recovery steps are:

\begin{enumerate}
  \item Run \code{poem check} to isolate language or include errors quickly.
  \item Run \code{poem doctor} if PDF compilation tools or fonts appear missing.
  \item Ensure every \directive{align} has an \directive{end}, and every include
        stays inside the project and ends in \code{.poem}.
  \item Run \code{poem build} again. If TeX itself fails, PoemTeX summarizes the
        useful diagnostic instead of exposing the full TeX stack trace.
\end{enumerate}

The compact normative specification is \code{FORMAT.md}; this guide is the
task-oriented companion. The project is an evolving open-source prototype, so
keep source files in version control and consult \code{CHANGELOG.md} when
upgrading.

\end{document}
